\documentclass[a4paper, 12pt, table, hidelinks, DIV=calc]{extarticle}

% ---------------------- Шрифты ----------------------
\usepackage{fontspec}
\setsansfont[
    Ligatures=TeX,
    BoldFont={Arial Narrow Bold},
    ItalicFont={Arial Narrow Italic},
    BoldItalicFont={Arial Narrow Bold Italic}
]{Arial Narrow}
\renewcommand\familydefault{\sfdefault}
% \setmainfont{Arial Narrow}

\usepackage{unicode-math}
\setmathfont{Stix Two Math}

% ---------------------- Язык, графика, цвета ----------------------
\usepackage[english,russian]{babel}
\usepackage{graphicx}
\graphicspath{{./images/}}

\usepackage{xcolor}
\definecolor{unimain}{RGB}{21, 73, 140}

\definecolor{unired}{RGB}{239, 134, 55} 
\definecolor{uniblue}{RGB}{21, 73, 140}
\definecolor{unigray}{RGB}{100, 100, 100}      
\definecolor{uniblack}{RGB}{20, 20, 20}
\definecolor{uniwhite}{RGB}{255, 255, 255} 
\definecolor{unigreen}{RGB}{11, 100, 54} %{0, 203, 98}
\definecolor{unidarkgreen}{RGB}{11, 100, 54}

% ---------------------- Геометрия ----------------------
\usepackage{typearea}
\usepackage[
    a4paper,
    includehead,
    includefoot,
    nomarginpar,
    top=6mm,
    bottom=6mm,
    left=22mm,
    right=10mm,
    headheight=14.5pt,
    headsep=3mm,
    footskip=7mm
]{geometry}
\savegeometry{standard}

% ---------------------- Абзацы ----------------------
\usepackage{parskip}
\setlength{\parindent}{1cm}
\setlength{\parskip}{1pt}
\setlength{\intextsep}{0mm}
\usepackage{indentfirst}

% ---------------------- Колонтитулы ----------------------
\usepackage{fancyhdr}

\fancypagestyle{plain}{%
  \fancyhf{}%
  \renewcommand{\headrulewidth}{1.5pt}%
  \renewcommand{\footrulewidth}{1.5pt}%
  \renewcommand{\headrule}{{\color{uniblue}\hrule height \headrulewidth}}%
  \renewcommand{\footrule}{{\color{uniblue}\hrule height \footrulewidth}}%
  \lhead{\color{black}\devicecode}%
  \chead{\color{unigray}Редакция \rev~от~\revdate}%
  \rhead{\color{black}\devicetype}%
  \lfoot{\color{black}Руководство по эксплуатации}%
  \rfoot{\thepage}%
}

% ---------------------- Таблицы ----------------------
\usepackage{longtable, booktabs}
\usepackage{calc}
\setlength{\LTpre}{5pt}
\setlength{\LTpost}{5pt}

\usepackage{colortbl}
\usepackage{tabularx}          % только один раз
\usepackage{multirow}
\usepackage{tabularray}        % если используется
\usepackage{hhline}
\usepackage{makecell}

% ---------------------- Подписи и нумерация ----------------------
\usepackage[labelsep=endash]{caption}
\DeclareCaptionLabelFormat{myformat}{\textls[75]{#1 #2}}
\captionsetup[figure]{name=Рисунок, justification=centering, labelformat=myformat} 
\captionsetup[table]{justification=justified, labelformat=myformat}
\setlength{\belowcaptionskip}{0mm}

% Нумерация рисунков и таблиц вида 1.1
\counterwithin*{figure}{section}
\counterwithin*{table}{section}
\renewcommand{\thefigure}{\thesection.\arabic{figure}}
\renewcommand{\thetable}{\thesection.\arabic{table}}

% ---------------------- Математика и формулы ----------------------
\usepackage{amsmath}
\usepackage{eqexpl}
\eqexplSetDelim{--}
\eqexplSetIntro{где}

% ---------------------- Заголовки ----------------------
\usepackage{titlesec}
\usepackage{appendix}
\titlespacing{\section}{1cm}{3pt}{3pt}
\titlespacing{\subsection}{1cm}{3pt}{3pt}
\titlespacing{\subsubsection}{1cm}{3pt}{3pt}

% Настройка расстояния между номером и текстом для section, subsection, subsubsection
\titleformat{\section}
  {\normalfont\Large\bfseries}
  {\thesection}
  {0.2cm}
  {}

% \section* (без номера) – по центру
\titleformat{name=\section, numberless}
  {\normalfont\Large\bfseries\centering}   % центрирование здесь
  {}          % номер не выводится
  {0pt}       % расстояние до текста (можно 0)
  {}          % код перед текстом (не используется)

\titleformat{\subsection}
  {\normalfont\large\bfseries}
  {\thesubsection}
  {0.2cm}
  {}

\titleformat{\subsubsection}
  {\normalfont\normalsize\bfseries}
  {\thesubsubsection}
  {0.2cm}
  {}

% ---------------------- Списки (enumitem) ----------------------
\usepackage{enumitem}
\setlist[itemize]{
    label=-,
    labelwidth=-0.5em,
    labelsep=0.5em,
    align=left,
    leftmargin=0em,
    itemindent=2.5em,
    nosep
}

\newlist{enumsec}{enumerate}{1}
\setlist[enumsec,1]{
    label=\arabic{section}.\arabic*,
    labelsep=-5pt,
    leftmargin=0pt,
    itemindent=47.5pt,
    align=left,
    nosep
}

\newlist{enumsub}{enumerate}{1}
\setlist[enumsub,1]{
    label=\arabic{section}.\arabic{subsection}.\arabic*,
    labelsep=4pt,
    leftmargin=0pt,
    itemindent=56.5pt,
    align=left,
    nosep      % убирает лишние вертикальные отступы между пунктами
}

\newlist{enumsubsub}{enumerate}{1}
\setlist[enumsubsub,1]{
    label=\arabic{section}.\arabic{subsection}.\arabic{subsubsection}.\arabic*,
    labelsep=5.5pt,
    leftmargin=0pt,
    itemindent=58pt,
    parsep=0pt,   % отступ между абзацами внутри одного пункта
    align=left,
    nosep
}

% ---------------------- Плавающие объекты ----------------------
\usepackage{float}
\usepackage{needspace}
\usepackage{placeins}

% ---------------------- Оглавление ----------------------
\usepackage{tocloft}
\usepackage{tocvsec2}   % оставлен, но используйте \setcounter{tocdepth} вместо \settocdepth
\renewcommand{\cftsecleader}{\cftdotfill{\cftdotsep}}
\renewcommand{\cftsecpagefont}{\normalfont}
\usepackage{etoolbox}
\makeatletter
\patchcmd{\l@section}{\cftsecfont #1}{\cftsecfont {#1}}{}{}
\makeatother
\renewcommand\cftsecfont{\MakeUppercase}

% ---------------------- Гиперссылки ----------------------
\usepackage{hyperref}

% ---------------------- Прочие пакеты ----------------------
\usepackage{pdflscape}
\usepackage{pdfpages}
\usepackage[absolute]{textpos}
\usepackage{fontawesome7}
\usepackage{microtype}
\usepackage{ulem}
\normalem
\usepackage{eso-pic}
\usepackage{subcaption}
\captionsetup[subfigure]{labelformat=empty}
\renewcommand{\thesubfigure}{\asbuk{subfigure})}
\makeatletter
\renewcommand{\p@subfigure}{\thefigure,\ }
\renewcommand{\thesubfigure}{\asbuk{subfigure}}
\makeatother
\usepackage{xstring}

% ---------------------- Сноски ----------------------
\usepackage[perpage]{footmisc}
\renewcommand{\thefootnote}{\arabic{footnote})}
\setlength{\footnotesep}{10pt}

% ---------------------- Прочее ----------------------
\usepackage{tikz}          % если используется
\uchyph=0                  % отключаем перенос в словах с заглавной

% ---------------------- Общее для РЭ ----------------------
\def\OKPD{Код ОКПД-2: 27.12.31.000} % ТУ 27.12.31-304-61775353-2025
\def\TNVED{Код ТН ВЭД: 8537 10 980 0 } % ТУ 27.12.31-304-61775353-2025
\def\ManualUnitGeneral{ЮТКБ.656122.611 РЭ «Устройства защиты и автоматики серии ЮНИТ-М500»}
\def\ManualTechSpec{ТУ 27.12.31-304-61775353-2025 «Низ\-ко\-вольт\-ные комплектные устройства систем защиты и автоматики серии ШЭТ»}
\def\twodoorcab{ЮТКБ.656463.300-001 РЭ, ЮТКБ.656463.300-002 РЭ, ЮТКБ.656463.300-003 РЭ, ЮТКБ.656463.300-004 РЭ,
 ЮТКБ.656463.300-031 РЭ, ЮТКБ.656463.300-032 РЭ, ЮТКБ.656463.300-033 РЭ, ЮТКБ.656463.300-051 РЭ}